% 翻译覆盖始于冻结源 XXVI.tex 第 2664 行；独立尾部，供主线程合并。
% BEGIN inputs/XXVI/tex/components/08_section5_1.tex
\SGASection{5}{共轭定理}
\SGARefTarget{sga3:xxvi:section:5}

\begin{sgastatement}{定理 5.1}
设 \(S\) 是半局部概形，\(G\) 是约化 \(S\)-群，\(P,P'\) 是 \(G\) 的两个相反抛物子群 \textup{（\SGARefLink{sga3:xxvi:definition:4-3-3}{4.3.3}）}。则
\[
 \rad^u(P)(S)\cdot P'\cdot P=G.
\]
换言之，当 \(u\) 遍历 \(\rad^u(P)(S)\) 时，\(G\) 的开子集 \(uP'\cdot P\) \textup{（\SGARefLink{sga3:xxvi:remark:4:3:6}{4.3.6}）} 覆盖整个 \(G\)。
\end{sgastatement}
\SGARefTarget{sga3:xxvi:theorem:5:1}

证明分为若干步。

\SGAInlineRefTarget{sga3:xxvi:paragraph:5-1-1}\paragraph{\SGARefLink{sga3:xxvi:paragraph:5-1-1}{5.1.1}.}
\SGARefTarget{sga3:xxvi:item:5:1:1}
只须在 \(S\) 为域 \(k\) 的谱时证明定理；这立即由 \SGARefLink{sga3:xxvi:cor:2:6}{2.6} 得出。

\SGAInlineRefTarget{sga3:xxvi:paragraph:5-1-2}\paragraph{\SGARefLink{sga3:xxvi:paragraph:5-1-2}{5.1.2}.}
\SGARefTarget{sga3:xxvi:item:5:1:2}
置 \(L=P\cap P'\)。假设 \(L\) 有 Borel 子群 \(B_L\)，并如第 XXII 讲 \SGARefLink{sga3:xxii:cor:5.9.7}{5.9.7}，取 \(B_L\) 的极大环面 \(T\)。立即可验
\[
 B=B_L\cdot\rad^u(P)
\]
是 \(P\) 的 Borel 子群。设 \(B'\) 是 \(G\) 中相对于 \(T\) 与 \(B\) 相反的 Borel 子群，因而 \(B\cap B'=T\)。使 \(G\) 相对于 \(T\) 分裂后，立即看出 \(B'\subset P'\)。我们断言
\[
 B^u(S)\cdot B'\cdot B
 \subset \rad^u(P)(S)\cdot P'\cdot P.                 \tag{x}\label{sga3:xxvi:theorem:5-1:equation:x}
\]
事实上，
\[
 B^u(S)\subset P(S)
 =\rad^u(P)(S)\cdot L(S)
 \subset\rad^u(P)(S)\cdot P'(S),
\]
故只须注意 \(B'\cdot B\subset P'\cdot P\)。于是 \textup{(x)} 将\SGARefLink{sga3:xxvi:theorem:5:1}{定理 5.1} 归结为有序对 \((B,B')\) 的情形。

\SGAInlineRefTarget{sga3:xxvi:paragraph:5-1-3}\paragraph{\SGARefLink{sga3:xxvi:paragraph:5-1-3}{5.1.3}.}
\SGARefTarget{sga3:xxvi:item:5:1:3}
当 \(k\) 代数闭时，定理成立：此时满足 \SGARefLink{sga3:xxvi:paragraph:5-1-2}{5.1.2} 的假设，可应用第 XXII 讲 \SGARefLink{sga3:xxii:prop:5.7.10}{5.7.10}。

\SGAInlineRefTarget{sga3:xxvi:paragraph:5-1-4}\paragraph{\SGARefLink{sga3:xxvi:paragraph:5-1-4}{5.1.4}.}
\SGARefTarget{sga3:xxvi:item:5:1:4}
当 \(k\) 是无限域时，定理也成立。事实上，由 \SGARefLink{sga3:xxvi:cor:2:7}{2.7}，
\[
 \rad^u(P)(k)\ \text{在}\ \rad^u(P)(\bar k)\ \text{中稠密},
\]
而定理在 \(\bar k\) 上成立。

\paragraph{5.1.5.}
\SGARefTarget{sga3:xxvi:item:5:1:5}
因而归结为 \(k\) 是有限域的情形。概形 \(\Bor(L)\) 是 \(L\) 下的光滑齐性空间；于是 Lang 定理 \textup{（Amer.\ J. Math.\ \textbf{78} \textup{(1956)}）} 表明 \(L\) 有 Borel 子群 \(B_L\)。
\footnote{\emph{编者注。} 另见 \textup{[DG70]}，\(\S\)III.5，\SGARefLink{sga3:xxvi:theorem:7-4}{7.4}。}
因此，由 5.1.2，可设 \(P=B\)、\(P'=B'\) 是 Borel 子群。记 \(T=B\cap B'\)。

\SGAInlineRefTarget{sga3:xxvi:paragraph:5-1-6}\paragraph{5.1.6.}
\SGARefTarget{sga3:xxvi:item:5:1:6}
设 \(K\) 是 \(k\) 的代数闭包。选取 \((G_K,B_K,T_K)\) 的钉扎，并以 \(R^+\)、\(\Delta\) 分别表示正根集、单根集。由第 XXII 讲 \SGARefLink{sga3:xxii:def:5.7.2}{5.7.2}，只须对每个 \(\alpha\in\Delta\) 证明
\[
 u_\alpha B'^{\,u}(K)
 \subset B^u(k)B'^{\,u}(K)B(K).                       \tag{1}\label{sga3:xxvi:theorem:5-1:equation:1}
\]
设 \(\alpha_i\) 是 \(\alpha\) 在 \(k\) 上的各个不同共轭元。由于 \(B\) 定义在 \(k\) 上，它们都属于 \(\Delta\)。设 \(R'\) 是可写为 \(\alpha_i\) 的线性组合的根的集合，并置 \(R'_- = R_-\cap R'\)。由于 \(R'\) 定义在 \(k\) 上，存在 \(T\) 的子环面 \(Q\)，使得 \(Q_K\) 是各 \(\alpha_i\) 的公共核的极大环面。

置
\[
 Z=\Centr_G(Q),\qquad B_Z=B\cap Z,\qquad B'_Z=B'\cap Z
\]
\textup{（第 XXII 讲 \SGARefLink{sga3:xxii:cor:5.10.2}{5.10.2}）}。只须在 \(Z\) 中验证所需断言，即
\[
 u_\alpha B_Z'^{\,u}(K)
 \subset B_Z^u(k)B_Z'^{\,u}(K)B_Z(K).                 \tag{2}\label{sga3:xxvi:theorem:5-1:equation:2}
\]
事实上，
\[
 (B_Z'^{\,u})_K=\prod_{\beta\in R'_-}U_\beta.
\]
设 \(R''=R-R'\)，并置
\[
 V=\prod_{\beta\in R''\cap R_-}U_\beta.
\]
则
\[
 (B'^{\,u})_K=(B_Z'^{\,u})_K\cdot V,
\]
且由第 XXII 讲 \SGARefLink{sga3:xxii:cor:5.6.7}{5.6.7}，\((B_Z)_K\) 正规化 \(V\)。因此，由 \textup{\SGARefLink{sga3:xxvi:theorem:5-1:equation:2}{(2)}} 依次得到
\[
\begin{aligned}
 u_\alpha B'^{\,u}(K)
 &=u_\alpha B_Z'^{\,u}(K)V(K)\\
 &\subset B_Z^u(k)B_Z'^{\,u}(K)B_Z(K)V(K)\\
 &\subset B_Z^u(k)B_Z'^{\,u}(K)V(K)B_Z(K),
\end{aligned}
\]
这立即给出 \textup{\SGARefLink{sga3:xxvi:theorem:5-1:equation:1}{(1)}}。于是归结为 \(G=Z\)，即 \(\operatorname{Gal}(K/k)\) 在单根集上传递作用的情形。

\SGAInlineRefTarget{sga3:xxvi:paragraph:5-1-7}\paragraph{\SGARefLink{sga3:xxvi:paragraph:5-1-7}{5.1.7}.}
\SGARefTarget{sga3:xxvi:item:5:1:7}
待证断言等价地说，\(G/B\) 是 \(B'^{\,u}\) 的开像在 \(B^u(k)\) 作用下各平移的并。将 \(G\) 换为其伴随群，甚至换为任意具有相同伴随群的群，都不改变这一断言。因此可设 \(G\) 是伴随群。

\SGAInlineRefTarget{sga3:xxvi:paragraph:5-1-8}\paragraph{\SGARefLink{sga3:xxvi:paragraph:5-1-8}{5.1.8}.}
\SGARefTarget{sga3:xxvi:item:5:1:8}
考虑 \(G_K\) 的 Dynkin 图。Galois 群在其上传递作用。但此 Galois 群只有循环商群，而 Dynkin 图没有回路。因此，该图为 \(nA_1\) 型，其中 \(n\geq0\)，或为 \(mA_2\) 型，其中 \(m\geq0\)。利用第 XXIV 讲 \SGARefLink{sga3:xxiv:proposition:5-9}{5.9} 的典范分解，可写成
\[
 G=\prod_{D/k}G_0,
\]
其中 \(D\) 是有限 \(k\)-概形，而 \(G_0\) 或为环面，或为 \(A_1\) 型、\(A_2\) 型。

由第 XXIV 讲 \SGARefLink{sga3:xxiv:corollary:5-12}{5.12}，群 \(B\) 来自 \(G_0\) 的 Borel 子群 \(B_0\)，而 \(T\) 来自 \(G_0\) 的极大环面 \(T_0\)；群 \(B'\) 来自 \(G_0\) 中相对于 \(T_0\) 与 \(B_0\) 相反的 Borel 子群 \(B'_0\)。有
\[
 B^{u}(k)=B_0^u(D),\qquad
 B'^{\,u}\cdot T\cdot B^u
 =\prod_{D/k}B_0'^{\,u}\cdot T_0\cdot B_0^u.
\]
因此，只须对三元组 \((G_0,B_0,T_0)\) 证明断言。

\SGAInlineRefTarget{sga3:xxvi:paragraph:5-1-9}\paragraph{\SGARefLink{sga3:xxvi:paragraph:5-1-9}{5.1.9}.}
\SGARefTarget{sga3:xxvi:item:5:1:9}
所以可设 \(G\) 为 \(\varnothing\)、\(A_1\) 或 \(A_2\) 型。由于 \(G\) 有 Borel 子群 \(B\)，由第 XXIV 讲 \SGARefLink{sga3:xxiv:3.9.1}{3.9.1}，它相对于 \(B\) 拟分裂，故在 \(\varnothing\) 型或 \(A_1\) 型时分裂。在分裂情形，已由第 XXII 讲 \SGARefLink{sga3:xxii:prop:5.7.10}{5.7.10} 知道定理成立，因此只剩下 \(A_2\) 型。

由第 XXIV 讲 \SGARefLink{sga3:xxiv:3.11}{3.11}，存在态射 \(E\to\operatorname{Spec}(k)\)，它是 \(A_2\) 型 Dynkin 图的自同构群 \(\mathbb Z/2\mathbb Z\) 下的 Galois 主齐性丛，使得
\[
 G=\operatorname{Q\text{-}\acute{E}p}_{E/\operatorname{Spec}(k)}
 (A_2).
\]
若 \(E\) 有截面，则 \(G\) 分裂，定理随即成立。否则 \(E=\operatorname{Spec}(k')\)，其中 \(k'/k\) 是二次扩张。最后，如 \SGARefLink{sga3:xxvi:paragraph:5-1-7}{5.1.7} 所述，可设 \(G\) 单连通，即 \(G\) 是 \(\operatorname{SL}_{3,k}\) 的形式。

\SGAInlineRefTarget{sga3:xxvi:paragraph:5-1-10}\paragraph{\SGARefLink{sga3:xxvi:paragraph:5-1-10}{5.1.10}.}
\SGARefTarget{sga3:xxvi:item:5:1:10}
现在处于如下情形：\(k\) 是有限域，\(k'/k\) 是二次扩张。对行列式为 \(1\) 的 \(3\times3\) 矩阵群 \(\operatorname{SL}_{3,k'}\)，如下选取钉扎：极大环面取对角矩阵群，Borel 子群取上三角矩阵群，根向量取为
\[
 u_\alpha=
 \begin{pmatrix}1&1&0\\0&1&0\\0&0&1\end{pmatrix},
 \qquad
 u_\beta=
 \begin{pmatrix}1&0&0\\0&1&1\\0&0&1\end{pmatrix}.
\]
大胞腔 \(\Omega\) 由下式给出：
\[
 \begin{pmatrix}a&b&c\\d&e&f\\g&h&i\end{pmatrix}
 \in\Omega(S)
 \quad\Longleftrightarrow\quad
 a\ \text{及}\ ae-bd\ \text{可逆},
\]
且
\[
 B^u(k)=
 \left\{
 \left.
 \begin{pmatrix}1&\bar x&z\\0&1&x\\0&0&1\end{pmatrix}
 \ \right|\
 x,z\in k',\ z+\bar z=x\bar x
 \right\}.
\]

须证明 \SGARefLink{sga3:xxvi:paragraph:5-1-6}{5.1.6} 中的包含关系 \textup{\SGARefLink{sga3:xxvi:theorem:5-1:equation:1}{(1)}}。因此，对任意 \(a,b,c\in K\)，须求 \(x,z\in k'\)，使得 \(z+\bar z=x\bar x\) 且
\[
 \begin{pmatrix}1&\bar x&z\\0&1&x\\0&0&1\end{pmatrix}
 \begin{pmatrix}1&1&0\\0&1&0\\0&0&1\end{pmatrix}
 \begin{pmatrix}1&0&0\\a&1&0\\b&c&1\end{pmatrix}
 \in\Omega(K).                                         \tag{3}\label{sga3:xxvi:theorem:5-1:equation:3}
\]
若 \(a\ne-1\)，取 \(x=z=0\)。若 \(a=-1\)，所需条件为
\[
 \begin{cases}
 z+\bar z=x\bar x,\\
 bz-\bar x\ne0,\\
 (b+c)\bar z+bx-1\ne0.
 \end{cases}                                           \tag{4}\label{sga3:xxvi:theorem:5-1:equation:4}
\]
设 \(q\geq2\) 是 \(k\) 的元素个数。对每个 \(m\in k\)，方程 \(z+\bar z=m\) 在 \(z\in k'\) 中有 \(q\) 个解。

若 \(b=0\)，取 \(x=1\)。只须解
\[
 z+\bar z=1,\qquad c\bar z\ne1,
\]
由上述观察，这总是可能的。

若 \(b\ne0\)，取 \(x=0\)。只须解
\[
 z+\bar z=0,\qquad z\ne0,\qquad (b+c)z\ne1.
\]
当 \(q\geq3\) 时，这总是可能的。当 \(k=\mathbb F_2\) 时，若 \(b+c\ne1\)，可取 \(z=1\)。

只剩下 \(k=\mathbb F_2\)、\(b+c=1\)、\(b\ne0\) 的情形。方程组变为
\[
 z+\bar z=x\bar x,\qquad bz\ne\bar x,\qquad
 \bar z+bx\ne1.
\]
若 \(b=1\)，或 \(b\notin k'\)，取 \(x=1\)；后两个条件变为
\[
 z\ne b^{-1},\qquad z\ne1-b,
\]
任意满足 \(z+\bar z=1\) 的解均满足这两个条件。最后，若 \(b\in k'-k\)，取 \(x=b\)、\(z=\bar b\)。证明完毕。
\SGAQED

% END inputs/XXVI/tex/components/08_section5_1.tex

% BEGIN inputs/XXVI/tex/components/09_section5_2_through_5_5.tex
\begin{sgastatement}{\SGAInlineRefTarget{sga3:xxvi:corollary:5-2}推论 5.2}
设 \(S\) 是半局部概形，\(G\) 是约化 \(S\)-群，\(P,P'\) 是 \(G\) 的两个相反抛物子群。典范映射
\[
 \rad^u(P)(S)\cdot\rad^u(P')(S)\longrightarrow(G/P)(S)
\]
是满射。特别地，
\[
 (G/P)(S)=G(S)/P(S).
\]
\(G\) 中每个与 \(P\) 同类型的抛物子群 \(Q\) 都形如
\[
 Q=\Int(uu')P,
\]
其中 \(u\in\rad^u(P)(S)\)、\(u'\in\rad^u(P')(S)\)。
\end{sgastatement}
\SGARefTarget{sga3:xxvi:cor:5:2}

第二个断言显然等价于第一个，故证明后者。设 \(s_i\) 是 \(S\) 的各闭点，\(V\) 是 \(G/P\) 中由 \(P'\) 的像给出的开子概形 \textup{（由 \SGARefLink{sga3:xxvi:remark:4:3:6}{4.3.6}，它同构于 \(\rad^u(P')\)）}，并取 \(x\in(G/P)(S)\)。由 \SGARefLink{sga3:xxvi:theorem:5:1}{5.1}，对每个 \(i\)，存在
\[
 u_i\in\rad^u(P)\bigl(\kappa(s_i)\bigr),
\]
使得 \(u_ix_{s_i}\) 是 \(V_{s_i}\) 的截面。由 \SGARefLink{sga3:xxvi:cor:2:6}{2.6}，可选取同时提升所有 \(u_i\) 的 \(u\in\rad^u(P)(S)\)。于是 \(ux\) 是 \(V\) 的截面，因为这可在闭纤维上检验。最后，
\[
 \rad^u(P')(S)\xrightarrow{\sim}V(S)
\]
是双射，结论得证。

\begin{sgastatement}{推论 5.3}
设 \(S\) 是半局部概形，\(G\) 是约化 \(S\)-群，\(P,P',Q\) 是 \(G\) 的三个抛物子群。存在 \(g\in G(S)\)，使得 \(\Int(g)Q\) 同时相对于 \(P\)、\(P'\) 处于横截位置。
\end{sgastatement}
\SGARefTarget{sga3:xxvi:cor:5:3}

采用 \SGARefLink{sga3:xxvi:corollary:4-2-4}{4.2.4}\textup{(ii)} 的记号，须证明 \(G\) 的普遍概形稠密开子概形
\[
 \Gen(Q/P)\cap\Gen(Q/P')
\]
在 \(S\) 上有截面。
\footnote{\emph{编者注。} Polo--Gille 版将“相对稠密”改为“在 \(S\) 上普遍概形稠密”；比较 EGA IV\textsubscript{3}，定义 11.10.8。}
由 \SGARefLink{sga3:xxvi:corollary:4-3-5}{4.3.5}\textup{(i)}，选取与 \(Q\) 相反的抛物子群 \(Q_1\)，并置
\[
 U=\rad^u(Q),\qquad U'=\rad^u(Q_1).
\]
我们将在 \(U(S)U'(S)\) 中求所需的 \(g\)。在这一形式下，\SGARefLink{sga3:xxvi:corollary:4-2-4}{4.2.4}\textup{(i)} 及 \SGARefLink{sga3:xxvi:cor:2:6}{2.6} 表明，只须在 \(S\) 的闭点上方的纤维上证明；因而可设 \(S=\operatorname{Spec}(k)\)。

若 \(k\) 代数闭，存在具有所需性质的 \(g\in G(k)\)。由 \SGARefLink{sga3:xxvi:corollary:5-2}{5.2}，可写成
\[
 g=uu'q,\qquad u\in U(k),\quad u'\in U'(k),\quad q\in Q(k),
\]
且 \(\Int(uu')Q=\Int(g)Q\)。

设 \(k\) 是无限域。令 \(V\) 是由笛卡尔图
\[
\begin{tikzcd}[column sep=huge,row sep=large]
 U\times_kU' \arrow[r] & G\\
 V \arrow[u,hook] \arrow[r]
   & \Gen(Q/P)\cap\Gen(Q/P') \arrow[u,hook].
\end{tikzcd}
\]
定义的 \(U\times_kU'\) 的开子集。由刚才所证，\(V(\bar k)\ne\varnothing\)，故 \(V\) 在 \(U\times_kU'\) 中稠密；于是由 \SGARefLink{sga3:xxvi:cor:2:7}{2.7}，它有 \(k\)-点。

最后设 \(k\) 是有限域。如 5.1.5 中那样应用 Lang 定理，可知 \(P\)、\(P'\) 分别有 Borel 子群 \(B\)、\(B'\)，因为
\[
 \Bor(P)\simeq\Bor(P/\rad^u(P)),\qquad
 \Bor(P')\simeq\Bor(P'/\rad^u(P'))
\]
是光滑的。
\footnote{\emph{编者注。} 法文底本的第二个同构中印作 \(\Bor(Q)\)；平行论证需要 \(\Bor(P')\)，此处采用后者。}
设 \(B_1\) 是与 \(B\) 相反的 Borel 子群。由 \SGARefLink{sga3:xxvi:corollary:5-2}{5.2}，存在
\[
 a\in B^u(k),\qquad a_1\in B_1^u(k),
\]
使得 \(\Int(aa_1)B=B'\)。于是
\[
 B_0=\Int(aa_1)B_1=\Int(a)B_1
\]
同时与 \(B'\)、\(B\) 相反。如 \SGARefLink{sga3:xxvi:lemma:3:8}{3.8}，设 \(Q_0\) 是 \(G\) 中包含 \(B_0\) 且与 \(Q\) 同类型的唯一抛物子群。由 \SGARefLink{sga3:xxvi:theorem:4-2-1}{4.2.1}\textup{\SGARefLink{sga3:xxvi:theorem:4-2-1:item:vi}{(vi)}}，它同时相对于 \(P\)、\(P'\) 处于横截位置。另一方面，\SGARefLink{sga3:xxvi:corollary:5-2}{5.2} 给出
\[
 Q_0=\Int(uu')Q,\qquad u\in U(k),\quad u'\in U'(k),
\]
从而断言得证。

\begin{sgastatement}{推论 5.4}
设 \(S\) 是半局部概形，\(G\) 是约化 \(S\)-群，\(P,Q\) 是 \(G\) 的两个抛物子群。存在 \(g\in G(S)\)，使得 \(\Int(g)P\) 相对于 \(Q\) 处于密切位置；等价地，\(\Int(g)P\cap Q\) 是 \(G\) 的抛物子群。
\end{sgastatement}
\SGARefTarget{sga3:xxvi:cor:5:4}

由 \SGARefLink{sga3:xxvi:corollary:4-3-5}{4.3.5}\textup{(i)}，选取与 \(P\) 相反的抛物子群 \(P'\)。由 \SGARefLink{sga3:xxvi:cor:5:3}{5.3}，存在与 \(P'\) 同类型、同时相对于 \(P\)、\(Q\) 处于横截位置的抛物子群 \(P'_1\)。如 \SGARefLink{sga3:xxvi:proposition:4-2-7}{4.2.7}\textup{(ii)}，设 \(T\) 是 \(P'_1\cap Q\) 的极大环面，并设 \(P_1\) 是相对于 \(T\) 与 \(P'_1\) 相反的抛物子群。由 \SGARefLink{sga3:xxvi:cor:4:4:5}{4.4.5}，\(P_1\)、\(Q\) 处于密切位置。由于 \(P\)、\(P_1\) 都与 \(P'_1\) 相反，由 \SGARefLink{sga3:xxvi:corollary:4-3-5}{4.3.5}\textup{(i)}，存在
\[
 g\in\rad^u(P'_1)(S),
\]
使得 \(\Int(g)P=P_1\)。
\SGAQED

同理，存在 \(u\in\rad^u(P)(S)\)，使得 \(\Int(u)P'=P'_1\)。故存在 \(u'\in\rad^u(P')(S)\)，使得 \(g=\Int(u)u'\)，且
\[
 P_1=\Int(uu'u^{-1})P=\Int(uu')P.
\]
这也给出了 \SGARefLink{sga3:xxvi:corollary:5-2}{5.2} 的另一个证明。

\SGARefLink{sga3:xxvi:cor:5:3}{推论 5.3} 及 \SGARefLink{sga3:xxvi:cor:5:4}{5.4} 是本节的主要结果。先记录 \SGARefLink{sga3:xxvi:cor:5:4}{5.4} 的若干推论。

\begin{sgastatement}{\SGAInlineRefTarget{sga3:xxvi:corollary:5-5}推论 5.5}
设 \(S\) 是半局部概形，\(G\) 是约化 \(S\)-群。
\begin{enumerate}[label=\textup{(\roman*)}]
\item 若 \(P,Q\) 是 \(G\) 的抛物子群，且 \(t(P)\subset t(Q)\)，则存在 \(g\in G(S)\)，使得 \(\Int(g)P\subset Q\)。
\item 设
\[
 P_1\supset P_2\supset\cdots\supset P_n,\qquad
 P'_1\supset P'_2\supset\cdots\supset P'_n
\]
是 \(G\) 的两条抛物子群链，且对每个 \(i\) 有 \(t(P_i)=t(P'_i)\)。存在 \(g\in G(S)\)，使得对每个 \(i\)，均有 \(\Int(g)P_i=P'_i\)。
\item 设 \(P,Q,P',Q'\) 是 \(G\) 的四个抛物子群，满足 \(t(P)=t(P')\)、\(t(Q)=t(Q')\)。若有序对 \((P,Q)\)、\((P',Q')\) 都处于横截位置，或都处于密切位置，则存在 \(g\in G(S)\)，使得
\[
 \Int(g)P=P',\qquad\Int(g)Q=Q'.
\]
\item 设 \(P,P'\) 是同类型的抛物子群，\(L\)、\(L'\) 分别是 \(P\)、\(P'\) 的 Levi 子群。存在 \(g\in G(S)\)，使得
\[
 \Int(g)P=P',\qquad\Int(g)L=L'.
\]
\end{enumerate}
\end{sgastatement}
\SGARefTarget{sga3:xxvi:cor:5:5}

\SGAProof
断言 \textup{(i)} 立即由 \SGARefLink{sga3:xxvi:cor:5:4}{5.4} 推出。断言 \textup{(ii)} 对 \(n\) 归纳证明，\(n=0\) 时显然。可设对 \(i=1,\ldots,n-1\)，有 \(P_i=P'_i\)。由 \SGARefLink{sga3:xxvi:corollary:5-2}{5.2}，存在 \(g\in G(S)\)，使得 \(\Int(g)P_n=P'_n\)。但 \(P_n\)、\(\Int(g)P_n\) 都包含于 \(P_{n-1}=P'_{n-1}\)，故由 \SGARefLink{sga3:xxvi:lemma:4:4:4}{4.4.4}\textup{(ii)}，
\[
 g\in P_{n-1}(S).
\]
从而对每个 \(i\)，均有 \(\Int(g)P_i=P'_i\)。

断言 \textup{(iv)} 立即由 \SGARefLink{sga3:xxvi:corollary:5-2}{5.2} 及 \SGARefLink{sga3:xxvi:cor:1:8}{1.8} 推出。先在横截情形证明 \textup{(iii)}。当 \(P\)、\(Q\) 的类型相反时，断言由 \textup{(iv)} 及 \SGARefLink{sga3:xxvi:definition:4-3-3}{4.3.3}\textup{(iii)} 推出。一般情形，由 \SGARefLink{sga3:xxvi:proposition:4-2-7}{4.2.7}\textup{(iii)} 及 \SGARefLink{sga3:xxvi:cor:4:4:6}{4.4.6}，可求得抛物子群
\[
 P_1\subset P,\quad P'_1\subset P',\quad
 Q_1\subset Q,\quad Q'_1\subset Q',
\]
使得 \(P_1,P'_1\) 相反，\(Q_1,Q'_1\) 相反，且 \(t(P_1)=t(P'_1)\)。于是存在 \(g\in G(S)\)，使得
\[
 \Int(g)P_1=P'_1,\qquad\Int(g)Q_1=Q'_1.
\]
施以共轭后，可设 \(P_1=P'_1\)、\(Q_1=Q'_1\)。此时 \(P,P'\) 处于密切位置且同类型，故由 \SGARefLink{sga3:xxvi:lemma:4:4:4}{4.4.4}\textup{(i)}，有 \(P=P'\)。同样，\(Q=Q'\)。

只剩下在密切情形证明 \textup{(iii)}。由 \SGARefLink{sga3:xxvi:corollary:5-2}{5.2}，可设 \(P=P'\)。再次由 \SGARefLink{sga3:xxvi:corollary:5-2}{5.2}，存在 \(g\in G(S)\)，使得 \(\Int(g)Q=Q'\)。但此时由 \SGARefLink{sga3:xxvi:lemma:4:4:4}{4.4.4}\textup{(ii)}，有 \(g\in P(S)\)，因此
\[
 \Int(g)P=P=P'.
\]
\SGAQED

% END inputs/XXVI/tex/components/09_section5_2_through_5_5.tex

% BEGIN inputs/XXVI/tex/components/10_section5_6_through_5_11.tex
\begin{sgastatement}{定义 5.6}
设 \(S\) 是概形，\(G\) 是约化 \(S\)-群，\(P\) 是 \(G\) 的抛物子群。若 \(G\) 中包含于 \(P\) 的每个抛物子群 \(Q\) 都等于 \(P\)，则称 \(P\) 是\emph{极小}的。
\end{sgastatement}
\SGARefTarget{sga3:xxvi:def:5:6}

注意，一般而言，此概念在取纤维时不保持。

\begin{sgastatement}{\SGAInlineRefTarget{sga3:xxvi:corollary:5-7}推论 5.7}
设 \(S\) 是半局部概形，\(G\) 是约化 \(S\)-群。
\begin{enumerate}[label=\textup{(\roman*)}]
\item 设
\[
 t,t'\in\Of\bigl(\operatorname{Dyn}(G)\bigr)(S).
\]
若 \(G\) 有一个类型为 \(t\) 的抛物子群及一个类型为 \(t'\) 的抛物子群，则它有类型为 \(t\cap t'\) 的抛物子群。特别地，当 \(P\) 遍历 \(G\) 的抛物子群时，所有 \(t(P)\) 组成的集合有最小元 \(t_{\min}\)。
\item \(G\) 的每个抛物子群都包含极小抛物子群。\(G\) 的抛物子群极小，当且仅当其类型为 \(t_{\min}\)。\(G\) 的任意两个极小抛物子群都由 \(G(S)\) 的某个元素共轭。
\end{enumerate}
\end{sgastatement}
\SGARefTarget{sga3:xxvi:cor:5:7}

这立即由 \SGARefLink{sga3:xxvi:cor:5:4}{5.4} 及 \SGARefLink{sga3:xxvi:corollary:5-5}{5.5}\textup{(i)} 推出。

\begin{sgastatement}{注 5.8}
与极小抛物子群相反的抛物子群也极小。因此，
\[
 s(t_{\min})=t_{\min}.
\]
\end{sgastatement}
\SGARefTarget{sga3:xxvi:rem:5:8}

\begin{sgastatement}{\SGAInlineRefTarget{sga3:xxvi:corollary:5-9}推论 5.9}
设 \(S\) 是概形，\(G\) 是约化 \(S\)-群，\(P\) 是 \(G\) 的抛物子群。典范态射
\[
 G\longrightarrow G/P=X
\]
使 \(G\) 成为 \(X\) 上的 Zariski 局部平凡 \(P_X\)-丛。若 \(L\) 是 \(P\) 的 Levi 子群，则典范态射 \textup{（参见 \SGARefLink{sga3:xxvi:cor:3:12}{3.12}）}
\[
 G\longrightarrow G/L=Y
\]
使 \(G\) 成为 \(Y\) 上的 Zariski 局部平凡 \(L_Y\)-丛。
\end{sgastatement}
\SGARefTarget{sga3:xxvi:cor:5:9}

只须证明如下提升断言：给定态射 \(S'\to S\)，其中 \(S'\) 为局部概形，每个态射 \(S'\to X\) \textup{（或 \(S'\to Y\)）} 都可提升为态射 \(S'\to G\)。换言之，可假设 \(S\) 是局部概形，而须证明
\[
 G(S)\longrightarrow X(S),
 \qquad\text{或分别为}\qquad
 G(S)\longrightarrow Y(S),
\]
是满射。

第一个断言已在 \SGARefLink{sga3:xxvi:corollary:5-2}{5.2} 中证明。对第二个断言，取 \(y\in Y(S)\)。它在 \(X(S)\) 中的典范像是某个 \(g\in G(S)\) 的像。因此，\(g\) 在 \(Y(S)\) 中的像 \(y'\) 与 \(y\) 在 \(X(S)\) 中的像相同。故存在唯一的
\[
 u\in\rad^u(P)(S),
\]
使得 \(y'u=y\)，而 \(gu\) 在 \(Y(S)\) 中的像确为 \(y\)。

\begin{sgastatement}{\SGAInlineRefTarget{sga3:xxvi:corollary:5-10}推论 5.10}
设 \(S\) 是半局部概形，\(G\) 是约化 \(S\)-群。
\begin{enumerate}[label=\textup{(\roman*)}]
\item 设 \(P\) 是 \(G\) 的抛物子群，\(L\) 是 \(P\) 的 Levi 子群。典范映射 \textup{（参见 3.21）} 诱导双射
\[
 H^1(S,L)\xrightarrow{\ \sim\ }H^1(S,P)
       \xrightarrow{\ \sim\ }H^1_{t(P)}(S,G).
\]
\item 若
\[
 t,t'\in\Of\bigl(\operatorname{Dyn}(G)\bigr)(S),
\]
则 \textup{（参见 3.21）}
\[
 H^1_t(S,G)\cap H^1_{t'}(S,G)=H^1_{t\cap t'}(S,G).
\]
\item 若 \(P\)、\(Q\) 是 \(G\) 中处于密切位置的两个抛物子群，则下列典范图是笛卡尔图，且其中所有箭头都是单射：
\[
\begin{tikzcd}[column sep=huge,row sep=large]
 H^1(S,P) \arrow[r] & H^1(S,G)\\
 H^1(S,P\cap Q) \arrow[u] \arrow[r] & H^1(S,Q). \arrow[u]
\end{tikzcd}
\]
\end{enumerate}
\end{sgastatement}
\SGARefTarget{sga3:xxvi:cor:5:10}

证明 \textup{(i)}。由 \SGARefLink{sga3:xxvi:cor:2:3}{2.3}，映射
\[
 H^1(S,L)\longrightarrow H^1(S,P)
\]
是双射，而由 3.21，
\[
 H^1(S,P)\longrightarrow H^1_{t(P)}(S,G)
\]
是满射。须证明它是单射，等价地，证明典范映射
\[
 H^1(S,P)\longrightarrow H^1(S,G)
\]
是单射。设 \(Q\) 是 \(P\) 下的主齐性丛，\(Q_1\) 是相伴的 \(G\) 下的主齐性丛，\(P'\)、\(G'\) 是相应的 \(P\)、\(G\) 的扭形式。显然，\(G'\) 是约化 \(S\)-群，\(P'\) 是 \(G'\) 的抛物子群。

\(H^1(S,P)\) 中在 \(H^1(S,G)\) 内与 \(Q\) 的类有相同像的元素，自然等同于
\[
 H^1(S,P')\longrightarrow H^1(S,G')
\]
的核。由上同调正合列，此核又等同于 \((G'/P')(S)\) 中的 \(G'(S)\)-轨道集。\footnote{关于非交换上同调中的这些论证，见 Giraud 的学位论文，\textup{[Gi71]}，\S~III.3。}
但由 \SGARefLink{sga3:xxvi:corollary:5-2}{5.2}，\(G'(S)\) 在 \((G'/P')(S)\) 上传递作用。

对 \textup{(ii)}，设 \(Q\) 是 \(G\) 下的主齐性丛，\(G_Q\) 是 \(G\) 的相应扭形式。由定义 \textup{（3.21）}，须证明 \(G_Q\) 有类型为 \(t\cap t'\) 的抛物子群，当且仅当它有类型为 \(t\)、\(t'\) 的抛物子群。这正是 \SGARefLink{sga3:xxvi:lemma:3:8}{3.8} 与 \SGARefLink{sga3:xxvi:corollary:5-7}{5.7}\textup{(i)} 的结合。最后，\textup{(iii)} 立即由 \textup{(i)}、\textup{(ii)} 推出。

现在叙述 \SGARefLink{sga3:xxvi:cor:5:3}{5.3} 的一个推论。

\begin{sgastatement}{推论 5.11}
设 \(S\) 是半局部概形，\(G\) 是约化 \(S\)-群，\(P\) 是 \(G\) 的抛物子群。若 \(P\ne G\)，则 \(G\) 至少有三个与 \(P\) 同类型的不同抛物子群。换言之，
\[
 P\ne G\qquad\Longrightarrow\qquad
 \bigl(G(S):P(S)\bigr)\geq 3.
\]
\end{sgastatement}
\SGARefTarget{sga3:xxvi:cor:5:11}

事实上，如 \SGARefLink{sga3:xxvi:corollary:4-3-5}{4.3.5}\textup{(i)}，设 \(P'\) 是与 \(P\) 相反的抛物子群。由于 \(P\ne G\)，有 \(\rad^u(P')\ne e\)，例如见 \SGARefLink{sga3:xxvi:item:4:3:1}{4.3.2}。由 \SGARefLink{sga3:xxvi:proposition:2-1}{2.1}，\(\rad^u(P')(S)\ne e\)；选取
\[
 u\in\rad^u(P')(S),\qquad u\ne e.
\]
则 \(\Int(u)P\ne P\)。由 \SGARefLink{sga3:xxvi:cor:5:3}{5.3}，存在与 \(P\) 同类型且同时与 \(P\)、\(\Int(u)P\) 相反的抛物子群 \(P_1\)。于是 \(P_1\)、\(P\)、\(\Int(u)P\) 是 \(G\) 的三个同类型的不同抛物子群。
\SGAQED

% END inputs/XXVI/tex/components/10_section5_6_through_5_11.tex

% BEGIN inputs/XXVI/tex/components/11_section6_1_through_6_8.tex
\SGASection{6}{抛物子群与分裂环面}
\SGARefTarget{sga3:xxvi:sec:6}

\begin{sgastatement}{\SGAInlineRefTarget{sga3:xxvi:proposition:6-1}命题 6.1}
设 \(S\) 是概形，\(G\) 是约化 \(S\)-群，\(\mathfrak g=\operatorname{Lie}(G/S)\)，\(Q\) 是 \(G\) 的分裂子环面。写成 \(Q=D_S(M)\)，并设
\[
 \mathfrak g=\bigoplus_{\alpha\in M}\mathfrak g^\alpha
\]
是 \(Q\) 作用下 \(\mathfrak g\) 的分解。设 \(M_1\subset M\) 满足
\[
 0\in M_1,\qquad
 \alpha,\beta\in M_1\Longrightarrow\alpha+\beta\in M_1.
\]
\begin{enumerate}[label=\textup{(\roman*)}]
\item 存在 \(G\) 的唯一光滑子群 \(H_{M_1}\)，它的纤维连通，包含 \(\Centr_G(Q)\)，且 Lie 代数为
\[
 \bigoplus_{\alpha\in M_1}\mathfrak g^\alpha.
\]
\item 有下列蕴涵：
\[
\begin{aligned}
 M_1=\{0\}
   &\Longrightarrow H_{M_1}=\Centr_G(Q),\\
 M_1=-M_1
   &\Longrightarrow H_{M_1}\text{ 是约化群},\\
 M_1\cup(-M_1)=M
   &\Longrightarrow
   \begin{gathered}
    H_{M_1}\text{ 与 }H_{-M_1}\text{ 是 \(G\) 的相反抛物子群，}\\
    \text{具有公共 Levi 子群 }H_{M_1\cap(-M_1)}.
   \end{gathered}
\end{aligned}
\]
\end{enumerate}
\end{sgastatement}
\SGARefTarget{sga3:xxvi:prop:6:1}

断言 \textup{(i)}、\textup{(ii)} 对 fpqc 拓扑是局部的。因此，可设 \(Q\) 包含于 \(G\) 的极大环面 \(T\)，并进一步使 \(G\) 相对于 \(T\) 分裂。此时，断言 \textup{(i)} 立即由\SGARefLink{sga3:xxii:cor:5.3.5}{第 XXII 讲 5.3.5}、\SGARefLink{sga3:xxii:prop:5.4.5}{5.4.5} 及 \SGARefLink{sga3:xxii:thm:5.4.7}{5.4.7} 推出。\textup{(ii)} 中的各断言由\SGARefLink{sga3:xxii:cor:5.3.5}{第 XXII 讲 5.3.5}、\SGARefLink{sga3:xxii:prop:5.10.1}{5.10.1}、\SGARefLink{sga3:xxii:lem:5.11.3}{5.11.3}，以及本讲的 \SGARefLink{sga3:xxvi:prop:1:3}{1.4}、\SGARefLink{sga3:xxvi:item:4:3:1}{4.3.2} 推出。

\begin{sgastatement}{推论 6.2}
设 \(S\) 是概形，\(G\) 是约化 \(S\)-群，\(Q\) 是 \(G\) 的分裂子环面。存在 \(G\) 的抛物子群，以 \(\Centr_G(Q)\) 为 Levi 子群。
\end{sgastatement}
\SGARefTarget{sga3:xxvi:cor:6:2}

事实上，写成 \(Q=D_S(M)\)，在群 \(M\) 上选取全序，并以 \(M_1\) 表示其正元素的集合。子群 \(H_{M_1}\) 具有所需性质。

\begin{sgastatement}{推论 6.3}
若约化 \(S\)-群 \(G\) 有非中心的分裂子环面，则它有真抛物子群 \textup{（即不同于 \(G\) 的抛物子群）}。
\end{sgastatement}
\SGARefTarget{sga3:xxvi:cor:6:3}

结合 \SGARefLink{sga3:xxvi:cor:6:2}{6.2}、\SGARefLink{sga3:xxvi:corollary:5-9}{5.9} 及 \SGARefLink{sga3:xxvi:corollary:5-10}{5.10}，得到如下结果。

\begin{sgastatement}{\SGAInlineRefTarget{sga3:xxvi:corollary:6-4}推论 6.4}
设 \(S\) 是概形，\(G\) 是约化 \(S\)-群，\(Q\) 是 \(G\) 的分裂子环面。典范态射
\[
 G\longrightarrow G/\Centr_G(Q)
\]
是局部平凡纤维化。

若 \(S\) 半局部，则映射
\[
 G(S)\longrightarrow\bigl(G/\Centr_G(Q)\bigr)(S)
\]
是满射，而映射
\[
 H^1\bigl(S,\Centr_G(Q)\bigr)\longrightarrow H^1(S,G)
\]
是单射。
\end{sgastatement}
\SGARefTarget{sga3:xxvi:cor:6:4}

\SGASubsection{6.5}{}
\SGARefTarget{sga3:xxvi:subsec:6:5}

设 \(S\) 连通。若 \(T\) 是 \(S\)-环面，而 \(T'\)、\(T''\) 是 \(T\) 的两个分裂子环面，则其乘积
\[
 T'\cdot T''
\]
仍是 \(T\) 的分裂子环面。\footnote{\emph{编者注。} Polo--Gille 版在此采用乘法记号，将底本的“和 \(T'+T''\)”改为“乘积 \(T'\cdot T''\)”。}
事实上，它等同于 \(T'\times_S T''\) 对 \(T'\cap T''\) 的商，由\SGARefLink{sga3:ix:proposition:2:11}{第 IX 讲 2.11}，此商分裂。从而 \(T\) 有最大的分裂子环面，记为 \(T_{\mathrm{spl}}\)。

\begin{sgastatement}{\SGAInlineRefTarget{sga3:xxvi:lemma:6-6}引理 6.6}
设 \(S\) 是连通概形，\(T\) 是等平凡 \(S\)-环面，\(T_{\mathrm{spl}}\) 是其最大的分裂子环面。下列条件等价：
\begin{enumerate}[label=\textup{(\roman*)}]
\item 存在不同于平凡同态的同态
\[
 T\longrightarrow\mathbf G_{m,S}.
\]
\item \(T_{\mathrm{spl}}\ne e\)。
\end{enumerate}
\end{sgastatement}
\SGARefTarget{sga3:xxvi:lem:6:6}

由于 \(T\) 等平凡，存在有限群 \(\Gamma\)、群为 \(\Gamma\) 的连通 Galois 主覆盖 \(S'\to S\)，以及同构
\[
 T_{S'}\simeq D_{S'}(M).
\]
此时群 \(M\) 是 \(\Gamma\)-模，且有自然同构
\[
 \operatorname{Hom}_S(T,\mathbf G_{m,S})
   =H^0(\Gamma,M).
\]

设 \(V\) 是由元素
\[
 \gamma(m)-m,\qquad \gamma\in\Gamma,\quad m\in M
\]
生成的 \(M\otimes\mathbf Q\) 的向量子空间。立即可验
\[
 (T_{\mathrm{spl}})_{S'}
   \simeq D_{S'}\bigl(M/(M\cap V)\bigr).
\]
因此，\textup{(i)} 等价于
\[
 H^0(\Gamma,M)\ne0,
\]
或等价地，
\[
 H^0(\Gamma,M\otimes\mathbf Q)\ne0,
\]
而 \textup{(ii)} 等价于 \(M\ne M\cap V\)，或等价于 \(M\otimes\mathbf Q\ne V\)。但
\[
 M\otimes\mathbf Q
 =H^0(\Gamma,M\otimes\mathbf Q)\oplus V.
\]
考虑从 \(M\otimes\mathbf Q\) 到其不变子空间的投影算子，它将 \(x\) 映为其在 \(\Gamma\) 下所有变换的平均值，即得此结论。

\begin{sgastatement}{引理 6.7}
设 \(S\) 是概形，\(G\) 是约化 \(S\)-群，\(P\ne G\) 是 \(G\) 的抛物子群，\(L\) 是 \(P\) 的 Levi 子群，\(Q=\rad(L)\)。存在不同于平凡同态的同态
\[
 Q\longrightarrow\mathbf G_{m,S}.
\]
\end{sgastatement}
\SGARefTarget{sga3:xxvi:lem:6:7}

设 \(U=\rad^u(P)\)。它在 \(\Int(P)\) 下不变，因而在 \(\Int(Q)\) 下不变。考虑可逆 \(\mathcal O_S\)-模
\[
 \det\bigl(\operatorname{Lie}(U)\bigr),
\]
即局部自由 \(\mathcal O_S\)-模 \(\operatorname{Lie}(U)\) 的最高外幂。伴随表示定义同态
\[
 f:Q\longrightarrow
 \operatorname{Aut}\!\left(\det\bigl(\operatorname{Lie}(U)\bigr)\right)
 =\mathbf G_{m,S}.
\]
由于 \(P\ne G\)，有 \(U\ne e\)。选取 \(s\in S\)，使得 \(U_s\ne e\)。相对于包含 \(Q_s\) 的极大环面使 \(G_s\) 分裂后，立即看出 \(f_s\ne e\)。

\begin{sgastatement}{命题 6.8}
设 \(S\) 是连通半局部概形，\(G\) 是约化 \(S\)-群，\(P\) 是 \(G\) 的抛物子群，\(L\) 是 \(P\) 的 Levi 子群，\(Q=\rad(L)\)。设 \(Q_{\mathrm{spl}}\) 是 \(Q\) 的最大分裂子环面，等价地，是 \(L\) 的最大中心分裂子环面。则
\[
 L=\Centr_G(Q_{\mathrm{spl}}).
\]
\end{sgastatement}
\SGARefTarget{sga3:xxvi:prop:6:8}

置 \(L'=\Centr_G(Q_{\mathrm{spl}})\)。这是 \(G\) 中包含 \(L\) 的约化子群，且由 \SGARefLink{sga3:xxvi:proposition:1-20}{1.20}，
\[
 P'=P\cap L'
\]
是 \(L'\) 的抛物子群，以 \(L\) 为 Levi 子群。若 \(L'\ne L\)，则 \(L'\not\subset P\)，因为由 \SGARefLink{sga3:xxvi:prop:1:6}{1.7}，\(L\) 是 \(P\) 的极大约化子群。故 \(P'\ne L'\)。

设 \(G_1\) 是 \(L'\) 的导出群，置 \(P_1=P'\cap G_1\)。由 \SGARefLink{sga3:xxvi:proposition:1-19}{1.19}，\(P_1\) 是半单群 \(G_1\) 的抛物子群，群 \(L_1=L\cap G_1\) 是 \(P_1\) 的 Levi 子群，且
\[
 Q_1=\rad(L_1)
    =\bigl(\rad(L)\cap G_1\bigr)^0
    =\bigl(Q\cap G_1\bigr)^0.
\]
由\SGARefLink{sga3:xiv:3.20}{第 XIV 讲 3.20}，群 \(L_1\) 有极大环面 \(T_1\)，而由\SGARefLink{sga3:xxiv:theorem:4-1-5}{第 XXIV 讲 4.1.5}，\(T_1\) 等平凡。因此，由\SGARefLink{sga3:ix:proposition:2:11}{第 IX 讲 2.11}，其子环面 \(Q_1\) 也等平凡。由于 \(P_1\ne G_1\)，\SGARefLink{sga3:xxvi:lem:6:7}{引理 6.7} 及 \SGARefLink{sga3:xxvi:lemma:6-6}{6.6} 给出
\[
 (Q_1)_{\mathrm{spl}}\ne e.
\]
从而
\[
 \bigl(Q_{\mathrm{spl}}\cap G_1\bigr)^0\ne e,
\]
因此 \(Q_{\mathrm{spl}}\not\subset\rad(L')\)，因为 \(\rad(L')\cap G_1\) 有限。这与 \(L'\) 的定义矛盾。
\SGAQED

% END inputs/XXVI/tex/components/11_section6_1_through_6_8.tex

% BEGIN inputs/XXVI/tex/components/12_section6_9_through_6_18.tex
\begin{sgastatement}{\SGAInlineRefTarget{sga3:xxvi:corollary:6-9}推论 6.9}
设 \(S\) 是连通半局部概形，\(G\) 是约化 \(S\)-群，\(Q\) 是 \(G\) 的临界子环面，即满足
\[
 \rad\bigl(\Centr_G(Q)\bigr)=Q
\]
的子环面。群 \(\Centr_G(Q)\) 是 \(G\) 的某个抛物子群的 Levi 子群，当且仅当
\[
 \Centr_G(Q)=\Centr_G(Q_{\mathrm{spl}}),
\]
或等价地，
\[
 \operatorname{Lie}(G)^Q
   =\operatorname{Lie}(G)^{Q_{\mathrm{spl}}}.
\]
\end{sgastatement}
\SGARefTarget{sga3:xxvi:cor:6:9}

这由 \SGARefLink{sga3:xxvi:cor:6:2}{6.2} 及 \SGARefLink{sga3:xxvi:prop:6:8}{6.8} 推出。

\begin{sgastatement}{\SGAInlineRefTarget{sga3:xxvi:corollary:6-10}推论 6.10}
设 \(S\) 是连通半局部概形，\(G\) 是约化 \(S\)-群，\(L\) 是 \(G\) 的子群。下列条件等价：
\begin{enumerate}[label=\textup{(\roman*)}]
\item 存在 \(G\) 的抛物子群，以 \(L\) 为 Levi 子群。
\item 存在 \(G\) 的分裂子环面，其中心化子为 \(L\)。
\item 存在同态
\[
 \mathbf G_{m,S}\longrightarrow G,
\]
其中心化子为 \(L\)。
\end{enumerate}
\end{sgastatement}
\SGARefTarget{sga3:xxvi:cor:6:10}

蕴涵 \textup{(i)}\(\Rightarrow\)\textup{(ii)} 由 \SGARefLink{sga3:xxvi:prop:6:8}{6.8} 推出，\textup{(iii)}\(\Rightarrow\)\textup{(i)} 由 \SGARefLink{sga3:xxvi:proposition:6-1}{6.1} 推出。只须证明 \textup{(ii)}\(\Rightarrow\)\textup{(iii)}。因此，假设
\[
 L=\Centr_G(Q),\qquad Q=D_S(M).
\]
记 \(\mathfrak g=\operatorname{Lie}(G/S)\)，并写成
\[
 \mathfrak g=\bigoplus_{\alpha\in M}\mathfrak g^\alpha,
\]
以 \(R\) 表示满足 \(\mathfrak g^\alpha\ne0\) 的 \(\alpha\in M\setminus\{0\}\) 的集合。

由于 \(R\) 有限且不含 \(0\)，存在同态
\[
 u:M\longrightarrow\mathbf Z,
\]
使得对每个 \(\alpha\in R\)，有 \(u(\alpha)\ne0\)。由对偶性，\(u\) 给出同态 \(\mathbf G_{m,S}\to Q\)，从而给出同态
\[
 f:\mathbf G_{m,S}\longrightarrow G.
\]
有
\[
 \Centr_G(f)\supset\Centr_G(Q).
\]
这是 \(G\) 的两个具有连通纤维的光滑子群，而且它们的 Lie 代数相同，因为都等于 \(\mathfrak g^0\)。因此，由通常的论证，它们相等。

\begin{sgastatement}{\SGAInlineRefTarget{sga3:xxvi:corollary:6-11}推论 6.11}
设 \(S\) 是连通半局部概形，\(G\) 是约化 \(S\)-群。映射
\[
 L\longmapsto\rad(L)_{\mathrm{spl}},
 \qquad
 Q\longmapsto\Centr_G(Q)
\]
在下列两个集合之间给出互逆的反序双射：
\begin{itemize}
\item \(G\) 的抛物子群的所有 Levi 子群 \(L\) 组成的集合；
\item 满足
\[
 \rad\bigl(\Centr_G(Q)\bigr)_{\mathrm{spl}}=Q
\]
的 \(G\) 的分裂子环面 \(Q\) 组成的集合。
\end{itemize}
\end{sgastatement}
\SGARefTarget{sga3:xxvi:cor:6:11}

\begin{sgastatement}{推论 6.12}
设 \(S\) 是连通半局部概形，\(G\) 是约化 \(S\)-群。考虑下列断言：
\begin{enumerate}[label=\textup{(\roman*)}]
\item \(G\) 有不同于 \(G\) 的抛物子群。
\item \(G\) 有非中心的分裂子环面。
\item[\textup{(ii bis)}] \(G\) 有相对维数为 \(1\) 的非中心分裂子环面。
\item 存在群同态
\[
 \mathbf G_{a,S}\longrightarrow G,
\]
它是闭浸入。
\end{enumerate}
则
\[
 \textup{(i)}
 \Longleftrightarrow\textup{(ii)}
 \Longleftrightarrow\textup{(ii bis)}
 \Longrightarrow\textup{(iii)}.
\]
\end{sgastatement}
\SGARefTarget{sga3:xxvi:cor:6:12}

唯一的新断言是 \textup{(i)}\(\Rightarrow\)\textup{(iii)}。设 \(P\ne G\) 是 \(G\) 的抛物子群。则 \(U=\rad^u(P)\ne e\)。如 \SGARefLink{sga3:xxvi:proposition:2-1}{2.1}，考虑 \(U\) 的合成列中最后一个非平凡子群 \(U_n\)。有同构
\[
 U_n\simeq W(E_n),
\]
其中 \(E_n\) 是局部自由 \(\mathcal O_S\)-模，由于 \(S\) 半局部且连通，它是自由模。由于 \(E_n\ne0\)，存在单态射
\[
 \mathcal O_S\longrightarrow E_n,
\]
其像局部地为直和项。它诱导闭浸入
\[
 \mathbf G_{a,S}=W(\mathcal O_S)
   \lhook\joinrel\longrightarrow W(E_n)\simeq U_n,
\]
从而 \textup{(iii)} 得证。

\begin{sgastatement}{注 6.12.1}
当 \(S\) 是特征为 \(0\) 的域的谱时，Jacobson--Morozov 定理给出 \textup{(iii)}\(\Rightarrow\)\textup{(ii bis)}。于是上述四个条件等价 \textup{（“Godement 判别法”；参见 \textup{[BT65]}，8.5）}。\footnote{更一般地，当 \(S\) 是完美域的谱时，这仍成立 \textup{（Tits）}；见 \textup{[BT71]}，推论 3.7。}
\end{sgastatement}
\SGARefTarget{sga3:xxvi:rem:6:12:1}

\begin{sgastatement}{定义 6.13}
设 \(S\) 是连通半局部概形，\(G\) 是约化 \(S\)-群。若 \(G\) 不含非平凡的分裂子环面，则称 \(G\) 是\emph{各向异性}的。
\end{sgastatement}
\SGARefTarget{sga3:xxvi:def:6:13}

\begin{sgastatement}{推论 6.14}
设 \(S\) 是连通半局部概形。约化 \(S\)-群 \(G\) 各向异性，当且仅当它没有抛物子群 \(P\ne G\)，且它的根基各向异性。
\end{sgastatement}
\SGARefTarget{sga3:xxvi:cor:6:14}

结合 \SGARefLink{sga3:xxvi:lemma:6-6}{6.6}、\SGARefLink{sga3:xxiv:theorem:4-1-5}{第 XXIV 讲 4.1.5} 及\SGARefLink{sga3:xxii:subsection:6-2}{第 XXII 讲 6.2}，得到：

\begin{sgastatement}{推论 6.15}
设 \(S\) 是连通半局部概形，\(G\) 是等平凡约化 \(S\)-群 \textup{（例如 \(G\) 半单，或 \(S\) 正规；参见\SGARefLink{sga3:x:theorem:5-16}{第 X 讲 5.16}）}。群 \(G\) 各向异性，当且仅当它没有抛物子群 \(P\ne G\)，且
\[
 \operatorname{Hom}_{S\text{-}\mathrm{gr}}
   (G,\mathbf G_{m,S})=e.
\]
\end{sgastatement}
\SGARefTarget{sga3:xxvi:cor:6:15}

\begin{sgastatement}{命题 6.16}
设 \(S\) 是连通半局部概形，\(G\) 是约化 \(S\)-群。\(G\) 的极大分裂子环面恰为 \(G\) 的极小抛物子群的 Levi 子群的最大中心分裂子环面。任意两个这样的环面都由 \(G(S)\) 的某个元素共轭。
\end{sgastatement}
\SGARefTarget{sga3:xxvi:prop:6:16}

设 \(Q\) 是 \(G\) 的极大分裂子环面。由 \SGARefLink{sga3:xxvi:cor:6:2}{6.2}，
\[
 L=\Centr_G(Q)
\]
是 \(G\) 的某个抛物子群 \(P\) 的 Levi 子群。由于
\[
 Q\subset\rad\bigl(\Centr_G(Q)\bigr)_{\mathrm{spl}},
\]
由 \(Q\) 的极大性得到
\[
 Q=\rad\bigl(\Centr_G(Q)\bigr)_{\mathrm{spl}}.
\]
由 \SGARefLink{sga3:xxvi:corollary:6-11}{6.11}，\(L\) 在 \(G\) 的抛物子群的 Levi 子群中极小。故由 \SGARefLink{sga3:xxvi:proposition:1-20}{1.20}，\(P\) 是 \(G\) 的极小抛物子群。由 \SGARefLink{sga3:xxvi:corollary:5-7}{5.7} 及 \SGARefLink{sga3:xxvi:corollary:5-5}{5.5}\textup{(iv)}，任意两个这样的子环面 \(Q\) 都由 \(G(S)\) 的某个元素共轭。环面 \(Q\) 及有序对 \((P,L)\) 的共轭性也证明了第一个断言。

\begin{sgastatement}{推论 6.17}
设 \(S\) 是连通半局部概形，\(P,P'\) 是处于标准位置的两个极小抛物子群 \textup{（\SGARefLink{sga3:xxvi:def:4:5:1:1}{4.5.1.1}）}。则 \(P\cap P'\) 包含 \(P\)、\(P'\) 的一个公共 Levi 子群。
\end{sgastatement}
\SGARefTarget{sga3:xxvi:cor:6:17}

事实上，由 \SGARefLink{sga3:xxvi:proposition:4-5-1}{4.5.1}\textup{(v)}，\(P\cap P'\) 包含 \(G\) 的极大环面 \(T\)。设 \(L\) 是 \(P\) 中包含 \(T\) 的唯一 Levi 子群。由 \SGARefLink{sga3:xxvi:proposition:1-21}{1.21}，
\[
 \rad(P)\cap T=\rad(P)\cap L=\rad(L).
\]
因此，\(\rad(P)\cap T\) 包含 \(\rad(L)_{\mathrm{spl}}\)，后者是 \(G\) 的极大分裂子环面，因而必等于 \(T_{\mathrm{spl}}\)。故
\[
 L=\Centr_G(T_{\mathrm{spl}}),
\]
由对称性，\(L\) 也是 \(P'\) 的 Levi 子群。

\begin{sgastatement}{注 6.18}
由 \SGARefLink{sga3:xxvi:proposition:1-21}{1.21} 可知，\(G\) 的抛物子群 \(P\) 极小，当且仅当 \(\rad(P)\) 包含 \(G\) 的极大分裂子环面。在此情形，若 \(T\) 是 \(P\) 的极大环面，则环面 \(T_{\mathrm{spl}}\) 同时是 \(G\) 及 \(\rad(P)\) 的极大分裂子环面，而 \(\Centr_G(T_{\mathrm{spl}})\) 是 \(P\) 的 Levi 子群。此外，\(P\) 的每个 Levi 子群都以此方式得到。
\end{sgastatement}
\SGARefTarget{sga3:xxvi:rem:6:18}

% END inputs/XXVI/tex/components/12_section6_9_through_6_18.tex

% BEGIN inputs/XXVI/tex/components/13_section7_opening_through_7_4.tex
\SGASection{7}{相对根数据}
\SGARefTarget{sga3:xxvi:sec:7}

本节始终以 \(S\) 表示非空连通半局部概形，\(G\) 表示约化 \(S\)-群，\(Q\) 表示 \(G\) 的极大分裂子环面，并以 \(L\) 表示 \(Q\) 在 \(G\) 中的中心化子：
\[
 L=\Centr_G(Q).
\]

\SGASubsection{7.1}{}
\SGARefTarget{sga3:xxvi:subsec:7:1}

由于 \(Q\) 是 \(L\) 的最大中心分裂子环面，\(G(S)\) 中每个正规化 \(L\) 的元素也正规化 \(Q\)。因此 \textup{（参见 \SGARefLink{sga3:xxvi:subsec:7:1:1}{7.1.1}）}
\[
 \Norm_G(L)(S)=\Norm_G(Q)(S).
\]
另一方面，\SGARefLink{sga3:xxvi:corollary:6-4}{6.4} 表明
\[
 G(S)\longrightarrow(G/L)(S)
\]
是满射。于是有典范等同
\[
 W_G(Q)(S)
 =\bigl(\Norm_G(Q)/\Centr_G(Q)\bigr)(S)
 \simeq\Norm_G(L)(S)/L(S).
\]

设
\[
 M=\operatorname{Hom}_{S\text{-}\mathrm{gr}}
      (Q,\mathbf G_{m,S}),
\]
故典范地有 \(Q\simeq D_S(M)\)。以 \(W\) 表示由 \(W_G(Q)(S)\) 定义的 \(M\) 的自同构群。于是
\[
 W\simeq W_G(Q)(S)\simeq\Norm_G(L)(S)/L(S).
\]

\SGASubsection{7.1.1}{}
\SGARefTarget{sga3:xxvi:subsec:7:1:1}

一般并没有 \(\Norm_G(L)=\Norm_G(Q)\)。例如，设 \(S\) 是域 \(k\) 的谱，该域有二次扩张 \(k'\)，并取 \(G\) 为酉群
\[
 {}^2\!A_{2,k'/k}
\]
\textup{（参见\SGARefLink{sga3:xxiv:3.11.2}{第 XXIV 讲 3.11.2}）}。\(G\) 的极小抛物子群就是其 Borel 子群，而它们的 Levi 子群是极大环面；故
\[
 \bigl(\Norm_G(L)/L\bigr)(\bar k)=\mathfrak S_3.
\]
\TranslationNote{英文底本在本段两条 Weyl 群等式的取点域中均遗漏了横线；经核对法文合卷 PDF 第 1361 页（印刷页 323、边页 504），此处均应在代数闭包上取点，故据法文补正。}
由于 \(G\) 不分裂，其极大分裂环面的维数至多为 \(1\)，因而同构于 \(\mathbf G_{m,k}\)。群 \(\Norm_G(Q)/L\) 忠实地作用在 \(Q\) 上，并且
\[
 \bigl(\Norm_G(Q)/L\bigr)(\bar k)=\mathbf Z/2\mathbf Z.
\]
\footnote{\emph{编者注。} 法文版在最后这个等式中印作 \(\Norm_G(Q)/Q\)，但紧接此前的作用以及相对 Weyl 群的定义均使用 \(\Norm_G(Q)/L\)。}

\SGASubsection{7.2}{}
\SGARefTarget{sga3:xxvi:subsec:7:2}

设 \(P\) 是 \(G\) 的一个以 \(L\) 为 Levi 子群的抛物子群；由 \SGARefLink{sga3:xxvi:cor:6:2}{6.2}，这样的子群存在。由 \SGARefLink{sga3:xxvi:rem:6:18}{6.18}，它必极小。利用 \(G\) 的极小抛物子群的共轭性 \textup{（\SGARefLink{sga3:xxvi:corollary:5-7}{5.7}）}、一个抛物子群的 Levi 子群的共轭性 \textup{（\SGARefLink{sga3:xxvi:cor:1:8}{1.8}）}，以及 \SGARefLink{sga3:xxvi:proposition:1-6}{1.6} 中的等式
\[
 P=\Norm_G(P),\qquad \Norm_G(L)\cap P=L,
\]
得到如下结论：\(G\) 中以 \(L\) 为 Levi 子群的 \textup{（极小）} 抛物子群组成的集合，是 \(W\) 下的主齐性空间。

\SGASubsection{7.3}{}
\SGARefTarget{sga3:xxvi:subsec:7:3}

\(G\) 的 Lie 代数在 \(Q\) 作用下分解为
\[
 \operatorname{Lie}(G)
 =\operatorname{Lie}(L)
  \oplus\bigoplus_{\alpha\in R}\operatorname{Lie}(G)^\alpha,
\]
其中 \(R\) 是 \(Q\) 的满足 \(\operatorname{Lie}(G)^\alpha\ne0\) 的非零特征 \(\alpha\) 组成的集合。这些就是 \(G\) 相对于 \(Q\) 的根。

设
\[
 M^*=\operatorname{Hom}_{S\text{-}\mathrm{gr}}
       (\mathbf G_{m,S},Q).
\]
它与 \(M\) 对偶，而 \(W\) 通过结构搬运自然地作用在它上面。

\begin{sgastatement}{\SGAInlineRefTarget{sga3:xxvi:theorem:7-4}定理 7.4}
采用 \SGARefLink{sga3:xxvi:subsec:7:3}{7.3} 的记号，存在唯一映射
\[
 R\longrightarrow M^*,\qquad \alpha\longmapsto\alpha^*,
\]
使 \((M,M^*)\) 成为 Weyl 群为 \(W\) 的根数据 \textup{（\SGARefLink{sga3:xxi:subsection:1-1}{第 XXI 讲 1.1}）}。

此外，\(G\) 中以 \(L\) 为 Levi 子群的抛物子群 \(P\)，通过关系
\[
 \operatorname{Lie}(P)
 =\operatorname{Lie}(L)
  \oplus\bigoplus_{\alpha\in R^+}\operatorname{Lie}(G)^\alpha,
\]
与 \(R\) 中的正根系 \(R^+\) 一一对应。
\end{sgastatement}
\SGARefTarget{sga3:xxvi:thm:7:4}

\SGASubsection{7.4.1}{}
\SGARefTarget{sga3:xxvi:subsec:7:4:1}

先假设已证明所需映射 \(\alpha\mapsto\alpha^*\) 的存在性。由\SGARefLink{sga3:XXI:3.4.10}{第 XXI 讲 3.4.10}，\(s_\alpha\) 是 \(W\) 中唯一使
\[
 s_\alpha(m)-m
\]
对每个 \(m\in M\) 都是 \(\alpha\) 的有理倍数的元素。因此，\(s_\alpha\) 由 \(\alpha\) 确定。由于
\[
 \alpha^*(m)\alpha=m-s_\alpha(m),
\]
\footnote{\emph{编者注。} 在此以及 \SGARefLink{sga3:xxvi:subsec:7:4:4}{7.4.4} 中，Polo--Gille 版均补回原文在 \(\alpha^*(m)\) 后遗漏的因子 \(\alpha\)。}
余根 \(\alpha^*\) 也由 \(\alpha\) 确定，从而证明了该映射的唯一性。

\SGASubsection{7.4.2}{}
\SGARefTarget{sga3:xxvi:subsec:7:4:2}

取 \(\alpha\in R\)。分别以 \(L_\alpha\)、\(H_\alpha\)、\(H_{-\alpha}\) 表示 \(G\) 中包含 \(L\)、纤维连通且满足下式的唯一光滑子群 \textup{（参见 \SGARefLink{sga3:xxvi:proposition:6-1}{6.1}\textup{(i)}）}：
\[
\begin{aligned}
 \operatorname{Lie}(L_\alpha)
  &=\operatorname{Lie}(L)
    \oplus\bigoplus_{\gamma\in\mathbf Z\alpha\cap R}
       \operatorname{Lie}(G)^\gamma,\\
 \operatorname{Lie}(H_\alpha)
  &=\operatorname{Lie}(L)
    \oplus\bigoplus_{\gamma\in\mathbf N\alpha\cap R}
       \operatorname{Lie}(G)^\gamma,\\
 \operatorname{Lie}(H_{-\alpha})
  &=\operatorname{Lie}(L)
    \oplus\bigoplus_{\gamma\in-\mathbf N\alpha\cap R}
       \operatorname{Lie}(G)^\gamma.
\end{aligned}
\]
由 \SGARefLink{sga3:xxvi:proposition:6-1}{6.1}\textup{(ii)}，群 \(L_\alpha\) 是 \(G\) 的约化子群，而 \(H_\alpha,H_{-\alpha}\) 是 \(L_\alpha\) 中以 \(L\) 为公共 Levi 子群的相反抛物子群。

因此，由 \SGARefLink{sga3:xxvi:subsec:7:2}{7.2}，存在元素
\[
 s_\alpha\in\Norm_{L_\alpha}(Q)(S)/L(S)\subset W,
\]
使得 \(s_\alpha(H_\alpha)=H_{-\alpha}\)。有
\[
 s_\alpha(\alpha)=-\alpha,
\]
因为 \(\alpha\)、\(-\alpha\) 分别是出现在 \(H_\alpha\)、\(H_{-\alpha}\) 中的 \(R\) 的元素的公因子。此外，\(s_\alpha^2=\mathrm{id}\)，因为 \(s_\alpha^2(H_\alpha)\)、\(H_\alpha\) 都相对于 \(L\) 与 \(H_{-\alpha}\) 相反。这样便构造了元素 \(s_\alpha\in W\)，具有性质
\[
\tag{x}\label{sga3:xxvi:subsection:7-4-2:equation:x}
 s_\alpha(\alpha)=-\alpha,\qquad s_\alpha^2=\mathrm{id},
\]
以及
\[
\tag{xx}\label{sga3:xxvi:subsection:7-4-2:equation:xx}
 s_\alpha\text{ 可由 }L_\alpha(S)\text{ 的一个元素代表}.
\]
该构造关于 \(\alpha\) 是典范的；特别地，
\[
\tag{xxx}\label{sga3:xxvi:subsection:7-4-2:equation:xxx}
 w s_\alpha w^{-1}=s_{w(\alpha)}
 \qquad\text{对每个 }w\in W.
\]

\SGASubsection{7.4.3}{}
\SGARefTarget{sga3:xxvi:subsec:7:4:3}

现在证明
\[
\tag{xxxx}\label{sga3:xxvi:subsection:7-4-3:equation:xxxx}
 s_\alpha(m)-m\in\mathbf Z\alpha
 \qquad\text{对每个 }m\in M.
\]
由于 \(S\) 连通，此断言对 fpqc 拓扑是局部的。因此可设
\[
 L_\alpha=G_1
\]
相对于 \(L\) 的极大环面 \(T_1\) 分裂。设 \((G_1,T_1,M_1,R_1)\) 是这样的分裂结构。单态射 \(Q\to T_1\) 将 \(M\) 等同于 \(M_1\) 的一个商；以
\[
 p:M_1\longrightarrow M
\]
表示典范映射。

\(R_1\) 在 \(p\) 下的像可能含零，由 \(G_1\) 相对于 \(Q\) 的根组成，因而由 \(R\) 中 \(\alpha\) 的整数倍组成。因此，
\[
 p(R_1)\subset\mathbf Z\alpha.
\]
由 \textup{\SGARefLink{sga3:xxvi:subsection:7-4-2:equation:xx}{(xx)}}，\(\Norm_{G_1}(L)(S)\) 的一个元素在 \(Q\) 上诱导 \(s_\alpha\)。于是\SGARefLink{sga3:xxii:rem:5.10.10}{第 XXII 讲 5.10.10} 给出截面
\[
 w\in(W_1)_S(S),
\]
它在 \(Q\) 上诱导 \(s_\alpha\)，其中 \(W_1\) 是根数据 \((M_1,R_1,\ldots)\) 的 Weyl 群。必要时缩小 \(S\)，可设 \(w\in W_1\)，且
\[
 p\bigl(w(m_1)\bigr)=s_\alpha\bigl(p(m_1)\bigr)
 \qquad(m_1\in M_1).
\]
由定义，\(w\) 是与 \(R_1\) 的元素相伴的反射的乘积。故 \(w(m_1)-m_1\) 是 \(R_1\) 的元素的整线性组合。从而
\[
 s_\alpha\bigl(p(m_1)\bigr)-p(m_1)
\]
是 \(p(R_1)\subset\mathbf Z\alpha\) 的元素的整线性组合，因此是 \(\alpha\) 的整数倍，正如所需。

\SGASubsection{7.4.4}{}
\SGARefTarget{sga3:xxvi:subsec:7:4:4}

现在可用
\[
 \alpha^*(m)\alpha=m-s_\alpha(m)
\]
定义 \(\alpha^*\in M^*\)。由 \textup{(x)}，
\[
 \langle\alpha^*,\alpha\rangle=2.
\]
此外，由 \textup{\SGARefLink{sga3:xxvi:subsection:7-4-2:equation:xxx}{(xxx)}}，对每个 \((\alpha,\alpha')\in R\times R\)，有
\[
 s_\alpha(\alpha')\in R,\qquad
 s_\alpha(\alpha'^*)=s_\alpha(\alpha')^*.
\]
因此，由\SGARefLink{sga3:xxi:subsection:1-1}{第 XXI 讲 1.1}，映射 \(\alpha\mapsto\alpha^*\) 在 \((M,M^*)\) 上定义了根数据。

\SGASubsection{7.4.5}{}
\SGARefTarget{sga3:xxvi:subsec:7:4:5}

设 \(W'\) 是此根数据的 Weyl 群，即由各 \(s_\alpha\) 生成的 \(M\) 的变换群。则 \(W'\subset W\)。

在自由交换群 \(M\) 上选取全序，并置
\[
 R^+=\{\alpha\in R\mid\alpha\geq0\}.
\]
这是 \(R\) 中的正根系。取 \(w\in W\)，并以元素
\[
 n\in\Norm_G(L)(S)=\Norm_G(Q)(S)
\]
代表它。采用 \SGARefLink{sga3:xxvi:proposition:6-1}{6.1} 的记号，置 \(P=H_{R^+}\)。则 \(P\) 是 \(G\) 中以 \(L\) 为 Levi 子群的抛物子群，且
\[
 \Int(n)P=H_{w(R^+)}.
\]
由 \SGARefLink{sga3:xxvi:subsec:7:3}{7.3}，\(w(R^+)=R^+\) 蕴涵 \(w=e\)。由\SGARefLink{sga3:XXI:3.3.7}{第 XXI 讲 3.3.7}，群 \(W'\) 在 \(R\) 的正根系集合上传递作用，而 \(W\) 中 \(R^+\) 的稳定子群平凡。因此，\(W=W'\)。

还可得出，\(W=W'\) 在 \(R\) 的正根系集合上以及 \(G\) 中以 \(L\) 为 Levi 子群的抛物子群集合上都单传递地作用。这证明了 \SGARefLink{sga3:xxvi:theorem:7-4}{7.4} 的最后一个断言。
\SGAQED

% END inputs/XXVI/tex/components/13_section7_opening_through_7_4.tex

% BEGIN inputs/XXVI/tex/components/14_section7_5_through_7_9.tex
\SGASubsection{7.5}{}
\SGARefTarget{sga3:xxvi:subsec:7:5}

设 \(P,P_1\) 是 \(G\) 的两个极小抛物子群，其 Levi 子群分别为 \(L,L_1\)，并设 \(Q,Q_1\) 分别是 \(L,L_1\) 的最大中心分裂子环面。则有序对 \((P,Q)\)、\((P_1,Q_1)\) 共轭：存在 \(g\in G(S)\)，使得
\[
 \Int(g)P=P_1,\qquad \Int(g)Q=Q_1.
\]
事实上，由 \SGARefLink{sga3:xxvi:corollary:5-7}{5.7}，\(P\)、\(P_1\) 共轭，故可设 \(P=P_1\)。此时由 \SGARefLink{sga3:xxvi:cor:1:8}{1.8}，\(L\)、\(L_1\) 由 \(P(S)\) 的一个元素共轭。此外，若 \(g,g'\) 是将 \((P,Q)\) 共轭到 \((P_1,Q_1)\) 的 \(G(S)\) 的两个元素，则 \(g'^{-1}g\) 正规化 \(P\)、\(Q\)，从而也正规化 \(P\)、\(L\)。但
\[
 \Norm_G(P)\cap\Norm_G(L)
 =P\cap\Norm_G(L)
 =L
 =\Centr_G(Q).
\]
因此，由 \(\Int(g)\) 诱导的同构
\[
 Q\xrightarrow{\ \sim\ }Q_1
\]
与 \(g\) 的选取无关。

设 \(\mathcal R,\mathcal R_1\) 是 \SGARefLink{sga3:xxvi:theorem:7-4}{7.4} 在
\[
 \operatorname{Hom}_{S\text{-}\mathrm{gr}}(Q,\mathbf G_{m,S})
 \quad\text{及}\quad
 \operatorname{Hom}_{S\text{-}\mathrm{gr}}(Q_1,\mathbf G_{m,S})
\]
中定义的根数据，并设 \(R^+,R_1^+\) 是对应于 \(P,P_1\) 的正根系。刚构造的典范同构 \(Q\simeq Q_1\) 将 \((\mathcal R,R^+)\) 搬运到 \((\mathcal R_1,R_1^+)\)。因此，利用上述具有传递性的同构系，将各个有序对 \((\mathcal R,R^+)\) 等同，便可定义 \(G\) 在 \(S\) 上的\emph{带基相对根数据}。\footnote{带基根数据是附带选定的正根系、等价地附带单根系的根数据；参见\SGARefLink{sga3:XXIII:1.5}{第 XXIII 讲 1.5}。}

今后，以
\[
 \mathcal R(G/S)
 =(\underline M,\underline M^*,\underline R,
   \underline R^*,\underline R^+)
\]
表示此带基根数据。因此，对每个上述有序对 \(P\supset Q\)，都有典范同构
\[
 \underline M\xrightarrow{\ \sim\ }
 \operatorname{Hom}_{S\text{-}\mathrm{gr}}(Q,\mathbf G_{m,S}),
\]
它将 \(\underline R\)、\(\underline R^+\) 分别搬运到 \(G\)、\(P\) 相对于 \(Q\) 的根，并将 \(W(\mathcal R)\) 搬运到 \(W_G(Q)(S)\)。

\SGASubsection{7.6}{}
\SGARefTarget{sga3:xxvi:subsec:7:6}

再设 \(Q\) 是 \(G\) 的极大分裂环面，\(P\) 是 \(G\) 中以 \(\Centr_G(Q)\) 为 Levi 子群的极小抛物子群，
\[
 (M,M^*,R,R^*,R^+)
\]
是 \SGARefLink{sga3:xxvi:theorem:7-4}{7.4} 中相应的带基根数据，\(\Delta\) 是 \(R^+\) 中的单根集。对每个 \(A\subset\Delta\)，置
\[
 R_A=R^+\cup(\mathbf Z A\cap R^-).
\]
因此，\(R_A\) 由正根以及可写成 \(A\) 的元素的整线性组合的负根组成。由\SGARefLink{sga3:XXI:3.1.4}{第 XXI 讲 3.1.4}，这是闭根集；而由\SGARefLink{sga3:XXI:3.3.10}{第 XXI 讲 3.3.10}，每个包含 \(R^+\) 的闭集都唯一地具有此形式。

由 \SGARefLink{sga3:xxvi:proposition:6-1}{6.1}，存在 \(G\) 的唯一光滑子群 \(P_A\)，它的纤维连通，包含 \(\Centr_G(Q)\)，并满足
\[
 \operatorname{Lie}(P_A)
 =\operatorname{Lie}(G)^0
  \oplus\bigoplus_{\alpha\in R_A}\operatorname{Lie}(G)^\alpha.
\]
由 \SGARefLink{sga3:xxvi:proposition:6-1}{6.1}、极小抛物子群的共轭性，以及 \(G\) 中包含 \(Q\) 的抛物子群的根集是闭集这一事实 \textup{（此事实在使群分裂后由 \SGARefLink{sga3:xxvi:prop:1:3}{1.4} 推出）}，立即得到下列结果。

\begin{sgastatement}{\SGAInlineRefTarget{sga3:xxvi:proposition:7-7}命题 7.7}
\begin{enumerate}[label=\textup{(\roman*)}]
\item 映射
\[
 A\longmapsto P_A
\]
给出从 \(\Delta\) 的子集集合到 \(G\) 中包含 \(P\) 的抛物子群集合的保包含关系的双射。
\item \(G\) 的每个抛物子群都由 \(G(S)\) 的某个元素共轭于唯一的 \(P_A\)。
\end{enumerate}
\end{sgastatement}
\SGARefTarget{sga3:xxvi:prop:7:7}

\SGASubsection{7.8}{}
\SGARefTarget{sga3:xxvi:subsec:7:8}

设 \(P\supset Q\) 如上。考虑 \SGARefLink{sga3:xxvi:subsec:7:5}{7.5} 中定义的 \(G\) 在 \(S\) 上的相对根数据及典范同构
\[
 f:(M,M^*,R,R^*,R^+)
   \xrightarrow{\ \sim\ }
   (\underline M,\underline M^*,\underline R,
    \underline R^*,\underline R^+).
\]
\(R^+\) 中的单根集 \(\Delta\) 被搬运到 \(\underline R^+\) 中的单根集 \(\underline\Delta\)。因此，每个子集 \(A\subset\Delta\) 都被搬运到子集
\[
 f(A)\subset\underline\Delta.
\]

\begin{sgastatement}{定义 7.9.0}
设 \(H\) 是 \(G\) 的任意抛物子群。由 \SGARefLink{sga3:xxvi:proposition:7-7}{7.7}\textup{(ii)}，它共轭于唯一的 \(P_A\)。置
\[
 t_r(H)=f(A)\subset\underline\Delta.
\]
由共轭定理立即可知，\(t_r(H)\) 与有序对 \(P\supset Q\) 的选取无关。称它为 \(H\) 的\emph{相对类型}。
\end{sgastatement}
\SGARefTarget{sga3:xxvi:def:7:9:0}

\begin{sgastatement}{\SGAInlineRefTarget{sga3:xxvi:proposition:7-9}命题 7.9}
\begin{enumerate}[label=\textup{(\roman*)}]
\item 映射 \(H\mapsto t_r(H)\) 诱导从 \(G\) 的抛物子群的 \(G(S)\)-共轭类集合到 \(\underline\Delta\) 的子集集合的双射。
\item 设 \(H\) 是 \(G\) 的抛物子群，\(P\) 是包含于 \(H\) 的极小抛物子群，\(Q\) 是 \(P\) 的某个 Levi 子群的最大中心分裂环面，\(\Delta\) 是 \(P\) 相对于 \(Q\) 的单根集，并设
\[
 f:\Delta\xrightarrow{\ \sim\ }\underline\Delta
\]
是典范同构。对每个 \(\alpha\in\Delta\)，
\[
 f(\alpha)\in t_r(H)
 \quad\Longleftrightarrow\quad
 \operatorname{Lie}(H)^{-\alpha}\ne0,
\]
且
\[
 H=P_A,\qquad A=f^{-1}\bigl(t_r(H)\bigr).
\]
\item 若 \(H,H'\) 是 \(G\) 中包含 \(P\) 的两个抛物子群，则 \textup{（其他等价条件见 \SGARefLink{sga3:xxvi:rem:3:8:1}{3.8.1}\textup{(ii)} 及 \SGARefLink{sga3:xxvi:corollary:5-5}{5.5}\textup{(i)}）}
\[
 t_r(H)\subset t_r(H')
 \quad\Longleftrightarrow\quad
 t(H)\subset t(H').
\]
\end{enumerate}
\end{sgastatement}
\SGARefTarget{sga3:xxvi:prop:7:9}

% END inputs/XXVI/tex/components/14_section7_5_through_7_9.tex

% BEGIN inputs/XXVI/tex/components/15_section7_10_and_7_11.tex
\SGASubsection{7.10}{}
\SGARefTarget{sga3:xxvi:subsec:7:10}

两个极小抛物子群的相对位置可如下描述。采用 4.5.2 的记号 \(t_2(P,P_1)\)。

\begin{enumerate}[label=\textup{(\arabic*)}]
\item 设 \(P,P_1,P',P'_1\) 是 \(G\) 的四个极小抛物子群。则
\[
 t_2(P,P_1)=t_2(P',P'_1),
\]
即有序对 \((P,P_1)\)、\((P',P'_1)\) 在 fpqc 拓扑下局部共轭，当且仅当存在 \(g\in G(S)\)，使得
\[
 \Int(g)P=P',\qquad \Int(g)P_1=P'_1.
\]

\item 固定 \(G\) 的一个极小抛物子群 \(P\)，以 \(L\) 为其 Levi 子群，并设 \(T\) 是 \(L\) 的极大环面。考虑 \(G\) 中相对于 \(P\) 处于标准位置的极小抛物子群的概形
\[
 \Par_{t_{\min}}(G;P).
\]
存在态射 \textup{（参见 \SGARefLink{sga3:xxvi:paragraph:4-5-5}{4.5.5}）}
\[
 f:\Par_{t_{\min}}(G;P)
   \longrightarrow W_P(T)\backslash W_G(T)/W_P(T),
\]
其纤维是 \(P\) 在 \(\Par_{t_{\min}}(G;P)\) 中的轨道。因此由 \textup{(1)}，\(f\) 诱导单态射
\[
 P(S)\backslash\Par_{t_{\min}}(G;P)(S)
 \lhook\joinrel\longrightarrow
 \bigl(W_P(T)\backslash W_G(T)/W_P(T)\bigr)(S).
\]
其像等同于 \(W\)。这就是 Bruhat 定理：\(P(S)\) 在 \(\Par_{t_{\min}}(G;P)(S)\) 中的每个轨道，恰含一个以 \(L\) 为 Levi 子群的 \(G\) 的抛物子群，即一个形如
\[
 \Int(n)P,\qquad n\in\Norm_G(L)(S)
\]
的子群。

\item 等价地，以 \(E(S)\) 表示使 \(\Int(g)P\) 与 \(P\) 处于相互标准位置的元素 \(g\in G(S)\) 的集合。则 \(E(S)\) 分拆为模 \(P(S)\) 的双陪集，以 \(W\) 为指标集：
\[
 E(S)=P(S)\,W\,P(S).
\]
置 \(U=\rad^u(P)\)，也可写成
\[
 E(S)=U(S)\cdot\Norm_G(L)(S)\cdot U(S).
\]
这给出 \(E(S)\) 按模 \(U(S)\) 的双陪集的分拆，以 \(\Norm_G(L)(S)\) 为指标集。

\item 若 \(S\) 是域的谱，则 \(E(S)=G(S)\)，从而恢复 \textup{[BT65]}，5.15。
\end{enumerate}

\begin{sgastatement}{\SGARefLink{sga3:xxvi:counterexamples:7:11}{反例 7.11}}
设
\[
 S=\operatorname{Spec}(\mathbf Z/4\mathbf Z),
 \qquad G=\operatorname{SL}_{2,S}.
\]
设 \(B\) 是通常的 Borel 子群，由满足
\[
 \begin{pmatrix}a&b\\c&d\end{pmatrix}
 \qquad\text{其中 }c=0
\]
的矩阵组成，并置
\[
 g=\begin{pmatrix}-1&1\\2&1\end{pmatrix}\in G(S),
 \qquad B'=\Int(g)B.
\]
则
\[
 B(S)=B'(S),
\]
但 \(B\cap B'\) 不含极大环面。
\footnote{事实上，\(B'\) 是 \(G\) 中由 \(c=2(a+b)\) 定义的子群。因此，\(B\cap B'\) 在 \(S\) 上不平坦，从而由 \SGARefLink{sga3:xxvi:proposition:4-5-1}{4.5.1}，不含极大环面。}

这同时表明，两个不同的极小抛物子群可以有相同的截面群，而且一般不能仅从群 \(P(S)\)、\(P'(S)\) 判断两个极小抛物子群 \(P,P'\) 是否处于标准位置。特别地，看来不能仅从配置
\[
 \{G(S),P(S),\Norm_G(L)(S)\}
\]
定义 \(G(S)\) 的子集 \(E(S)\)。在上例中，该子集由 \(c\ne2\) 定义。
\footnote{更一般地，对每个 \(S\)-概形 \(S'\)，集合 \(E(S')\) 由 \(c\) 为零或可逆的条件定义。}
\end{sgastatement}
\SGARefTarget{sga3:xxvi:counterexamples:7:11}

% END inputs/XXVI/tex/components/15_section7_10_and_7_11.tex

% BEGIN inputs/XXVI/tex/components/16_section7_12_through_bibliography.tex
\SGASubsection{7.12}{}
\SGARefTarget{sga3:xxvi:subsec:7:12}

现在研究 \(\mathcal R(G/S)\) 如何随 \(S\) 变化。设 \(S'\) 是同样非空、连通且半局部的 \(S\)-概形。若 \(Q\) 是 \(G\) 的极大分裂环面，则 \(Q_{S'}\) 是 \(G_{S'}\) 的分裂环面。选取 \(G_{S'}\) 中包含 \(Q_{S'}\) 的极大分裂环面 \(Q'\)，并置
\[
\begin{aligned}
 M&=\operatorname{Hom}_{S\text{-}\mathrm{gr}}
       (Q,\mathbf G_{m,S})
   \simeq
   \operatorname{Hom}_{S'\text{-}\mathrm{gr}}
       (Q_{S'},\mathbf G_{m,S'}),\\
 M'&=\operatorname{Hom}_{S'\text{-}\mathrm{gr}}
       (Q',\mathbf G_{m,S'}).
\end{aligned}
\]
单态射 \(Q_{S'}\to Q'\) 诱导满态射
\[
 u:M'\longrightarrow M.
\]
记
\[
 L=\Centr_G(Q),\qquad L'=\Centr_{G_{S'}}(Q').
\]
则 \(L'\subset L_{S'}\)。

若 \(H\) 是 \(G\) 中包含 \(L\) 的子群，则 \(H_{S'}\) 包含 \(L'\)，且
\[
\begin{aligned}
 \operatorname{Lie}(H)
  &=\operatorname{Lie}(L)
    \oplus\bigoplus_{\alpha\in R_H}
       \operatorname{Lie}(H)^\alpha,\\
 \operatorname{Lie}(H_{S'})
  &=\operatorname{Lie}(L')
    \oplus\bigoplus_{\alpha'\in R'_{H_{S'}}}
       \operatorname{Lie}(H_{S'})^{\alpha'}.
\end{aligned}
\]
这里 \(R_H\)、\(R'_{H_{S'}}\) 分别表示 \(H\)、\(H_{S'}\) 相对于 \(Q\)、\(Q'\) 的根集。立即得到
\[
 R_H\subset u\bigl(R'_{H_{S'}}\bigr)
       \subset R_H\cup\{0\}.
\]

取 \(H=G\)，首先得到
\[
 R\subset u(R')\subset R\cup\{0\}.
\]
再取 \(H\) 为以 \(L\) 为 Levi 子群的极小抛物子群 \(P\)。则 \(R'_{H_{S'}}\) 包含 \(R'\) 的一个正根系。因此，由 \SGARefLink{sga3:xxvi:theorem:7-4}{7.4}，存在 \(G_{S'}\) 的极小抛物子群 \(P'\)，以 \(L'\) 为 Levi 子群且包含于 \(P_{S'}\)。这样便构造了包含关系图
\[
\begin{array}{ccccc}
 Q_{S'} & \subset & L_{S'} & \subset & P_{S'}\\[1.2ex]
 \cap   &         & \cup   &         & \cup\\[1.2ex]
 Q'     & \subset & L'     & \subset & P'.
\end{array}
\]

若 \(R^+\)、\(\Delta\) 分别是 \(P\) 定义的 \(R\) 的正根系、单根集，并类似地定义 \(R'^+,\Delta'\)，则
\[
 R^+\subset u(R'^+)\subset R^+\cup\{0\},
 \qquad
 \Delta\subset u(\Delta')\subset\Delta\cup\{0\}.
\]

现在取
\[
 w\in W\simeq\Norm_G(Q)(S)/\Centr_G(Q)(S),
\]
并选取代表元 \(n\in\Norm_G(Q)(S)\)。由于 \(\Int(n)Q=Q\)，有 \(\Int(n)L=L\)，从而 \(\Int(n)L_{S'}=L_{S'}\)。因此，环面 \(Q'\)、\(\Int(n)Q'\) 是 \(L_{S'}\) 的两个极大分裂环面，故由某个 \(x\in L(S')\) 共轭。可选取 \(x\)，使得
\[
 \Int(nx)Q'=Q'.
\]
于是 \(n'=nx\) 属于 \(\Norm_{G_{S'}}(Q')(S')\)。设 \(w'\) 是 \(n'\) 在
\[
 W'\simeq
 \Norm_{G_{S'}}(Q')(S')/\Centr_{G_{S'}}(Q')(S')
\]
中的像。\(w'\) 在 \(M'\) 上的作用与 \(u:M'\to M\) 相容，且在 \(M\) 上诱导的作用就是 \(w\) 的作用。

利用相对根数据的定义以及共轭定理，得到如下结果。

\begin{sgastatement}{\SGAInlineRefTarget{sga3:xxvi:theorem:7-13}定理 7.13}
设 \(S,S'\) 是非空连通半局部概形，\(S'\to S\) 是概形态射，\(G\) 是约化 \(S\)-群，并以
\[
\begin{aligned}
 \mathcal R(G/S)
  &=(\underline M,\underline M^*,\underline R,
     \underline R^*,\underline R^+),\\
 \mathcal R(G_{S'}/S')
  &=(\underline M',\underline M'^*,\underline R',
     \underline R'^*,\underline R'^+)
\end{aligned}
\]
表示带基相对根数据。存在典范同态
\[
 u:\underline M'\longrightarrow\underline M,
\]
具有如下性质：
\begin{enumerate}[label=\textup{(\roman*)}]
\item 同态 \(u\) 是满射。
\item 对每个 \(w\in W\)，存在与 \(u\) 相容并在 \(\underline M\) 上诱导 \(w\) 的元素 \(w'\in W'\)。
\item 对每个子集 \(X\subset\underline M\)，置
\[
 X^\wedge=X\cap(\underline M\setminus\{0\}).
\]
则
\[
 u(\underline R'^+)^\wedge=\underline R^+,
 \qquad
 u(\underline\Delta')^\wedge=\underline\Delta.
\]
\item 对每个抛物子群 \(H\)，考虑
\[
 t_r(H)\subset\underline\Delta,
 \qquad
 t_r(H_{S'})\subset\underline\Delta'.
\]
则
\[
\begin{aligned}
 t_r(H_{S'})
 &=u^{-1}\bigl(t_r(H)\cup\{0\}\bigr)
       \cap\underline\Delta'\\
 &=\left\{\alpha'\in\underline\Delta'
       \ \middle|\ 
       u(\alpha')\in t_r(H)\text{ 或 }u(\alpha')=0
   \right\}.
\end{aligned}
\]
\end{enumerate}
\end{sgastatement}
\SGARefTarget{sga3:xxvi:thm:7:13}

\begin{sgastatement}{注 7.14}
若 \(G\) 在 \(S\) 上可分裂，则其极大分裂环面就是极大环面，这里引入的相对概念与此前定义的绝对概念一致。因此，前一定理利用 \(G_{S'}\) 的绝对根数据和绝对类型，描述相对根数据 \(\mathcal R(G/S)\) 及相对类型 \(t_r\)，其中选取 \(S'\) 使 \(G_{S'}\) 可分裂 \textup{（参见\SGARefLink{sga3:xxiv:theorem:4-4-1}{第 XXIV 讲 4.4.1}）}。这一描述见 \textup{[BT65]}，\SGARefLink{sga3:xxvi:cor:6:11}{6.12} 及随后各节。
\end{sgastatement}
\SGARefTarget{sga3:xxvi:rem:7:14}

\SGASubsection{7.15}{}
\SGARefTarget{sga3:xxvi:subsec:7:15}

设 \(S\) 是 Hensel 局部概形，\(s_0\) 是其闭点，\(S_0\) 是 \(s_0\) 处剩余域的谱，将其等同于 \(S\) 的闭子概形。对 \(S\) 上的每个对象 \(X\)，以 \(X_0\) 表示其到 \(S_0\) 的基变换。设 \(G\) 是约化 \(S\)-群。

对 \(G\) 的每个抛物子群 \(\overline P\)，群 \(\overline P_0\) 是 \(G_0\) 的抛物子群。反之，对 \(G_0\) 的每个抛物子群 \(P\)，存在 \(G\) 的抛物子群 \(\overline P\)，使得
\[
 \overline P_0=P.
\]
这由 Hensel 引理及 \(\Par(G)\) 是光滑 \(S\)-概形这一事实推出。特别地，由 \SGARefLink{sga3:xxvi:corollary:5-7}{5.7}，\(\overline P\) 极小，当且仅当 \(\overline P_0\) 极小。

选定这样的子群 \(\overline P\) 后，同样的论证表明，\(\overline P_0\) 的每个极大分裂子环面都形如 \(T_0\)，其中 \(T\) 是 \(\overline P\) 的极大分裂子环面。容易由此得出，\(G\) 在 \(S\) 上及 \(G_0\) 在 \(S_0\) 上的相对根数据典范同构。因此，\(G\) 的抛物子群理论归结为 \(G_0\) 的相应理论。

反之，由\SGARefLink{sga3:xxiv:1.21}{第 XXIV 讲命题 1.21}，每个约化 \(S_0\)-群都形如 \(G_0\)。这使得约化 \(S_0\)-群的抛物子群研究可归结为在 \(S\) 上的相应研究。

\chapter*{参考文献}
\addcontentsline{toc}{chapter}{参考文献}
\SGARefTarget{sga3:xxvi:bibliography}

\begin{description}[leftmargin=4.5em,style=nextline]
\item[\textup{[BT65]}]
A.~Borel 与 J.~Tits，\emph{Groupes réductifs}，
Publ.\ Math.\ I.H.É.S.\ \textbf{27} \textup{(1965)}，55--150。
\footnote{\emph{编者注。} 此文献见于原版；以下文献由 Polo--Gille 版增补。}

\item[\textup{[BT71]}]
A.~Borel 与 J.~Tits，\\
\emph{Éléments unipotents et sous-groupes
paraboliques de groupes réductifs. I}，
Invent.\ Math.\ \textbf{12} \textup{(1971)}，95--104。

\item[\textup{[Ch05]}]
C.~Chevalley，\emph{Classification des groupes algébriques
semi-simples}，P.~Cartier、A.~Grothendieck 与 M.~Lazard 合作，
文集第~3~卷，Springer，2005。

\item[\textup{[DG70]}]
M.~Demazure 与 P.~Gabriel，\emph{Groupes algébriques}，
Masson \& North-Holland，1970。

\item[\textup{[Gi71]}]
J.~Giraud，\emph{Cohomologie non abélienne}，
Springer-Verlag，1971。
\end{description}

% END inputs/XXVI/tex/components/16_section7_12_through_bibliography.tex
